% !TeX program = xelatex
\documentclass[11pt]{article}

\usepackage[a4paper, margin=1in]{geometry}
\usepackage{amsmath, amssymb, amsthm}
\usepackage{stmaryrd}
\usepackage{listings}
\usepackage{xcolor}
\usepackage{hyperref}
\usepackage{microtype}
\usepackage{fontspec}
\usepackage{booktabs}
\setmonofont{DejaVu Sans Mono}[Scale=0.85]

\sloppy
\setlength{\emergencystretch}{3em}

\newcommand{\btt}[1]{\texttt{\detokenize{#1}}}

% ---- Lean code listing style -------------------------------------------
\lstdefinelanguage{lean}{
  morekeywords={theorem, lemma, def, by, fun, let, in, do, match, with,
    if, then, else, return, import, open, namespace, end, structure,
    inductive, partial, private, noncomputable, syntax, elab, elab_rules,
    macro, macro_rules, tactic, conv, command, example, have, show, intro,
    apply, simp, rw, ring, all_goals, repeat, decide, native_decide,
    fin_cases, norm_num, exact, set_option},
  sensitive=true,
  morecomment=[l]{--},
  morecomment=[s]{/-}{-/},
  morestring=[b]"
}
\lstset{
  inputencoding=utf8,
  language=lean,
  basicstyle=\ttfamily\footnotesize,
  keywordstyle=\color{blue!70!black}\bfseries,
  commentstyle=\color{gray}\itshape,
  stringstyle=\color{red!60!black},
  showstringspaces=false,
  breaklines=true,
  columns=fullflexible,
  frame=single,
  framesep=4pt,
  rulecolor=\color{black!30},
  backgroundcolor=\color{gray!6},
  xleftmargin=6pt,
  xrightmargin=6pt,
  aboveskip=0.8\baselineskip,
  belowskip=0.8\baselineskip
}
% -----------------------------------------------------------------------

\title{What Can Be Generated Automatically}
\author{}
\date{}

\begin{document}
\maketitle

This note lists what the repository can currently generate without hand
transcription: the command to run, what it produces, and when you would
use it. It is a reference for the state of the tooling, not an argument
about it. Names follow the convention \texttt{\_n\_k\_m\_i}: $n$ vertices,
type $\sigma$ the $k$-vertex graph of index $m$ (\texttt{0\_0} is the empty
type $\emptyset_t$), enumeration index $i$.

One fact is worth stating up front because it affects how you use all of
the generators: every generated definition comes with a lemma closed by
\texttt{decide}/\texttt{native\_decide} or a verified tactic, and the
result is re-checked by the Lean kernel. So if a generator's input is
wrong, you find out as a \texttt{lake build} failure on the generated
file, not as a wrong theorem. You can therefore treat the scripts as
untrusted helpers and let the build decide whether the output is sound.

\section{The pieces and how they stack}

There are three groups of generators. Each consumes only what the group
below it produced.

\begin{center}
\begin{tabular}{lll}
\toprule
Group & Produces & Needs as input \\
\midrule
1. Flags        & flag/type definitions, completeness lemmas & nothing \\
2. Numerical    & density / multiplication theorems & density JSON (\S\ref{sec:pipeline}) \\
3. Certificate  & SDP matrices, vectors, the main theorem & a flagmatic JSON \\
\bottomrule
\end{tabular}
\end{center}

A typical use runs them top to bottom: the flag definitions and the
density data are built once and committed; from then on, porting a new
certificate is just Group~3.

\section{Group 1: flag and type definitions}
\label{sec:flags}

\paragraph{Where.} \texttt{Flags/FlagGenerator.lean} defines the macros;
\texttt{Flags/FlagDef.lean} invokes them, one line per size. These run at
elaboration time and read no files.

\paragraph{\texttt{generate\_empty\_typed\_flags n}.} Produces, for every
non-isomorphic $n$-vertex graph $i$: \texttt{Sym2Graph\_n\_0\_0\_i},
\texttt{Sym2Flag\_n\_0\_0\_i}, \texttt{Flag\_n\_0\_0\_i},
\texttt{FlagAlgebra\_n\_0\_0\_i}, plus the finsets and completeness lemmas
\texttt{flagSet\_n\_0\_0}, \texttt{flagSet\_n\_0\_0\_val\_eq},
\texttt{flagSet\_n\_0\_0\_eq\_univ} (and their \texttt{sym2} counterparts).
Currently instantiated for $n = 0,\dots,6$.

\paragraph{\texttt{generate\_flags k m n}.} Produces the type
\texttt{FlagType\_k\_m} and, for every flag $i$ on it:
\texttt{Sym2LabeledGraph\_n\_k\_m\_i}, \texttt{Sym2Flag\_n\_k\_m\_i},
\texttt{Flag\_n\_k\_m\_i}, \texttt{FlagAlgebra\_n\_k\_m\_i}, the
\texttt{@[simp]} lemmas \texttt{unlabel\_n\_k\_m\_i} and
\texttt{downward\_n\_k\_m\_i}, and the matching \texttt{flagSet\_n\_k\_m\_*}
lemmas. Currently instantiated for the type/size pairs the present
scenarios use, up to $n = 5$.

\paragraph{When to use.} Only when you need flags of a size or type not
already in \texttt{FlagDef.lean} -- add the corresponding
\texttt{generate\_*} line and rebuild. For everything else the constants
are already available project-wide, so you just refer to them by name.

\paragraph{Limits.} $n = 7$ (1044 graphs) is not instantiated: the
elaboration-time enumeration runs in the interpreter and takes over 80
minutes, and the completeness bridge's list literal exceeds the default
\texttt{maxRecDepth}.

\section{Group 2: density and multiplication theorems}
\label{sec:data}

\paragraph{Where.} \texttt{Flags/Densities/DensityLoader.lean} and
\texttt{Flags/Densities/MulLoader.lean}. Each macro reads one JSON file
(produced by \S\ref{sec:pipeline}) and registers \texttt{simp} theorems
about the Group-1 constants. Put the relevant lines near the top of your
proof file.

\paragraph{\texttt{load\_forbid\_density\_theorems "*\_free\_indices.json"}.}
For each $n$-vertex empty-typed flag, registers
$\mathrm{flagDensity}_1\,H\,\texttt{Flag\_n\_0\_0\_i} = 0$ when the flag is
$H$-free and $\ne 0$ otherwise. Recognizes the tags $K_3$, $K_4$, $K_5$.
Used to drop forbidden host flags.

\paragraph{\texttt{load\_flag\_pair\_density\_theorems "density\_*.json"}.}
Registers $\mathrm{flagDensity}_2\,\texttt{Flag\_<pat>\_p1}\,\texttt{Flag\_<pat>\_p2}\,\texttt{Flag\_<host>\_h} = \text{value}$,
one per row of the file. Used wherever a pair density appears as a
\texttt{simp} fact.

\paragraph{\texttt{load\_forbid\_mul\_theorems "density\_*.json"}.} From the
same file, registers the product expansions
$\texttt{FlagAlgebra\_<pat>\_i}\cdot\texttt{FlagAlgebra\_<pat>\_j}
=_H \sum_h c\,\texttt{FlagAlgebra\_<host>\_h}$
for each $H$-free pair $(i,j)$, named
\texttt{flagMul\_FlagAlgebra\_<pat>\_i\_FlagAlgebra\_<pat>\_j}. These are
what \texttt{reduce\_downward\_flagmul} looks up in the main proof. A
companion \texttt{load\_mul\_theorems} emits the plain (no-forbid) equality
version.

\paragraph{File-name discipline.} The loaders require the forbid to be in
the file name (\texttt{\_forbid\_<H>}, or \texttt{\_no\_forbid}). This is
deliberate: two density tables computed under different forbids would
otherwise share a name, and you could load the wrong one. Keep the tag in
the name and this cannot happen.

\section{The data pipeline behind Group 2}
\label{sec:pipeline}

Group~2 reads JSON; this is where that JSON comes from. Four Python
scripts under \texttt{Flags/} and \texttt{Flags/Densities/}, in a fixed
order:

\begin{lstlisting}[language={}]
generate_graphs.py
   └─→ Graphs/graphs_n.json          (canonical edge lists, n ≤ 7)
         ├─→ gen_free_indices.py
         │     └─→ Densities/graphs_n_<H>_free_indices.json
         │           ▷ load_forbid_density_theorems
         └─→ generate_flags.py  (also reads graphs_k.json for σ)
               └─→ Flags/flags_n_k_m.json
                     └─→ calculate_densities.py  (host + pattern)
                           └─→ Densities/density_<host>_from_<pat>_forbid_<H>.json
                                 ▷ load_flag_pair_density_theorems
                                 ▷ load_forbid_mul_theorems
\end{lstlisting}

These JSON files feed \emph{only} the Group-2 loaders. They used to also be
the source of the Group-1 flag definitions, but those are now built inside
Lean (\S\ref{sec:flags}); the scripts' docstrings still mention the old
\texttt{FlagLoader.lean} consumer, which is stale.

\paragraph{\texttt{generate\_graphs.py}.} Lists every non-isomorphic
$n$-vertex graph via NetworkX's atlas ($n \le 7$), canonicalizes each to
its lexicographically smallest edge list, and writes
\texttt{Graphs/graphs\_n.json}. The size $n$ is hard-coded in the
\texttt{\_\_main\_\_} block. Run it once per new $n$.

\paragraph{\texttt{generate\_flags.py n k type\_num}.} Reads
\texttt{graphs\_k.json} (the type $\sigma$) and \texttt{graphs\_n.json} (the
underlying graphs), enumerates flags by automorphism orbit, and writes
\texttt{Flags/flags\_n\_k\_m.json} with each flag's
\texttt{type\_indices} and \texttt{downward\_coeff}.
\begin{lstlisting}[language={}]
python generate_flags.py 4 2 0
\end{lstlisting}

\paragraph{\texttt{gen\_free\_indices.py <graphs\_n.json> {-}{-}forbid-H}.}
Writes the indices of the $H$-free graphs to
\texttt{graphs\_n\_<H>\_free\_indices.json}. $H$ is \texttt{--forbid-K3} /
\texttt{-K4} / \texttt{-K5}, or \texttt{--forbid <flagmatic string>}
(e.g.\ \texttt{3:122331}).
\begin{lstlisting}[language={}]
python gen_free_indices.py ../Graphs/graphs_4.json --forbid-K3
\end{lstlisting}

\paragraph{\texttt{calculate\_densities.py {-}{-}host \ldots {-}{-}pattern \ldots {-}{-}forbid-H}.}
Computes the exact pair density $p(F_1, F_2; G)$ over all $H$-free hosts
and forbidden-free pattern pairs, and writes
\texttt{density\_<host>\_from\_<pat>\_forbid\_<H>.json} (or
\texttt{\_no\_forbid} if \texttt{--forbid} is omitted).
\begin{lstlisting}[language={}]
python calculate_densities.py --host ../Flags/flags_4_2_0.json \
       --pattern ../Flags/flags_3_2_0.json --forbid-K3
\end{lstlisting}

\paragraph{On index ordering.} The scripts order graphs and flags by their
own canonical form; Lean (Group~1) orders them by its own. The density
JSON refers to flags only by index, so the two orderings have to match. If
they ever did not, the \texttt{native\_decide} on the generated density
lemma would fail -- the value would be asserted about the wrong
Lean-defined flag -- so a mismatch shows up at build time. You do not have
to check the ordering by hand.

\section{Group 3: porting a certificate to a Lean module}
\label{sec:cert}

\paragraph{Where.} \texttt{Flagmatic/flagmatic\_to\_lean.py}. Input is one
flagmatic SDP certificate JSON; output is a \texttt{.lean} module under
\texttt{Flagmatic/}. Subcommands:
\begin{lstlisting}[language={}]
python flagmatic_to_lean.py inspect      <cert>.json
python flagmatic_to_lean.py check-deps   <cert>.json
python flagmatic_to_lean.py gen-skeleton <cert>.json <target>.lean [--namespace N]
\end{lstlisting}

\paragraph{\texttt{inspect}.} Read-only. Prints the Lean identifier each
flagmatic string in the certificate resolves to. Run it first on a new
certificate to confirm the strings map to constants that exist.

\paragraph{\texttt{check-deps}.} Reports which density JSON files
(\S\ref{sec:pipeline}) the certificate will need, marks each present or
missing, and prints the \texttt{load\_*} lines that will be emitted. Exits
nonzero if anything is missing, so it doubles as a pre-build guard. If it
reports a missing file, generate it with the relevant pipeline script,
then re-run.

\paragraph{\texttt{gen-skeleton}.} Writes the full module: the SDP matrices
$M_t$ with their LDL$^\top$ positive-semidefiniteness proofs, the real
liftings, the \texttt{load\_*} lines, the type \texttt{σ} and flag vectors
\texttt{v}, the main theorem, and -- when the objective lives below host
size -- an \texttt{*\_expand\_under\_forbid} lemma plus its
\texttt{flagDensity}$_1$ evaluation table. The result builds with no
\texttt{sorry}.

\paragraph{Workflow.}
\begin{lstlisting}[language={}]
inspect  →  check-deps  →  gen-skeleton  →  lake build
\end{lstlisting}
The first two catch identifier and dependency problems before any Lean
compilation.

\section{Worked example (Mantel)}

The full sequence for $K_3$-free edge density at $N=3$, assuming the
$n\le 3$ flags are already in \texttt{FlagDef.lean}:
\begin{lstlisting}[language={}]
# data the certificate needs (run once, output committed)
python gen_free_indices.py ../Graphs/graphs_3.json --forbid-K3
python calculate_densities.py --host ../Flags/flags_3_1_0.json \
       --pattern ../Flags/flags_2_1_0.json --forbid-K3

# port the certificate
python flagmatic_to_lean.py check-deps   Certificates/mantel_cert.json
python flagmatic_to_lean.py gen-skeleton Certificates/mantel_cert.json \
       Mantel.lean --namespace Mantel
lake build
\end{lstlisting}
The generated \texttt{Mantel.lean} loads the two JSON files, states
$\texttt{FlagAlgebra\_2\_0\_0\_1} \le_{K_3} \tfrac12\cdot 1$, and proves it.

\section{Adding a new forbidden graph}
\label{sec:newforbid}

\emph{Update: for the complete graphs $K_n$ this five-file procedure has since
been collapsed to a single line --- see Section~\ref{sec:knupdate}. The recipe
below is the general one, still used as written for non-$K_n$ forbids
($C_4$, $P_4$, \ldots).}

To forbid a graph $X$ (on $n_X$ vertices) beyond $K_3$, $K_4$, $K_5$, you
touch five files, grouped into the three steps below. The certificate
translator (\S\ref{sec:cert}) is \emph{not} one of them: it finds $X$ by
re-parsing \texttt{CommonGraphs.lean}, so it needs no edit. Before
starting, find the canonical index $i_X$ of $X$ in
\texttt{graphs\_$n_X$.json} (e.g.\ with \texttt{inspect}), and make sure
the size-$n_X$ flags are already generated in \texttt{FlagDef.lean}.

\begin{enumerate}
\item \textbf{\texttt{Forbid/CommonGraphs.lean}} -- define $X$ and its flag
  identity:
\begin{lstlisting}
def X : SimpleGraph (Fin n_X) := ...
lemma X_toFinFlag_eq : X.toFinFlag = ⟨n_X, Flag_nX_0_0_iX⟩ := by ...
\end{lstlisting}
  Copy the $K_3$/$K_4$/$K_5$ entries as a template; only the \texttt{def}
  body and the index $i_X$ change.

\item \textbf{The two data scripts} (\texttt{gen\_free\_indices.py},
  \texttt{calculate\_densities.py}) -- add a \texttt{--forbid-X} shorthand.
  The subgraph-containment logic is already generic (a plain
  \texttt{--forbid <flagmatic string>} works for any graph); what the
  shorthand adds is a clean tag \texttt{"X"}, which is what lands in the
  JSON's \texttt{forbid.tag} field and in the output file name. That tag
  has to match the Lean identifier, which is why the raw \texttt{--forbid
  3:\ldots} form (tag \texttt{"3\_\ldots"}) is not enough by itself.

\item \textbf{The two forbid loaders} -- the step that is easy to miss.
  Both \texttt{load\_forbid\_density\_theorems}
  (\texttt{DensityLoader.lean}) and \texttt{load\_forbid\_mul\_theorems}
  (\texttt{MulLoader.lean}) \texttt{match} on the tag and currently know
  only \texttt{"K3"}, \texttt{"K4"}, \texttt{"K5"} -- any other tag throws
  \texttt{Unsupported forbid tag}. Add a \texttt{| "X" => \ldots} branch to
  each, copied from an existing one, replacing the tag, the lemma name
  \texttt{X\_toFinFlag\_eq}, and the index in
  \texttt{simp [Flag\_nX\_0\_0\_iX]}. Each branch is about 15 lines. (The
  pair-density loader \texttt{load\_flag\_pair\_density\_theorems} and the
  no-forbid \texttt{load\_mul\_theorems} are generic and need no branch.)
\end{enumerate}

Then generate the data and build:
\begin{lstlisting}[language={}]
python gen_free_indices.py ../Graphs/graphs_nX.json --forbid-X
python calculate_densities.py --host ... --pattern ... --forbid-X
python flagmatic_to_lean.py gen-skeleton <cert>.json <target>.lean
lake build
\end{lstlisting}

Step~3 is manual because each forbid lemma bakes \texttt{X\_toFinFlag\_eq}
into its proof script, so the loaders dispatch on the tag rather than
abstracting over the graph. At three forbids the copy-paste has been
cheaper than that abstraction; if the list grows it is the natural thing
to factor out.

\section{Update: generalizing the $K_n$ forbids}
\label{sec:knupdate}

The procedure of Section~\ref{sec:newforbid} treated every forbidden graph
alike, complete graphs included. But $K_n$ is the one forbid family that is
fully parametric in a single number, so it no longer needs that procedure. The
three places that used to hard-code $K_3$, $K_4$, $K_5$ one at a time have each
been collapsed to a single $n$-generic form.

\paragraph{Python --- was three shorthands, now one.}
\btt{gen_free_indices.py} and \btt{calculate_densities.py} carried separate
\btt{--forbid-K3}, \btt{--forbid-K4}, \btt{--forbid-K5} flags, each a
hand-written edge tuple. They now share one \btt{--forbid-Kn N}: it builds the
$K_N$ edge set from a one-line \btt{complete_graph_edges} helper and emits the
clean tag \btt{KN} for any $N$. The three old flags stay as aliases routing
through the same path, so existing commands and committed file names are
unchanged. (The raw \btt{--forbid <flagmatic>} form still handles arbitrary
graphs; what it lacks is the clean \btt{KN} tag the loaders key on --- which is
why the shorthands existed in the first place.)

\paragraph{\texttt{CommonGraphs.lean} --- was three def/lemma pairs, now three
one-liners.} The hand-written \btt{def K3} / \btt{K3_toFinFlag_eq} and its
$K_4$, $K_5$ copies are replaced by a macro \btt{generate_complete_graph r idx},
which emits \btt{def K{r}} and proves
\btt{K{r}_toFinFlag_eq : K{r}.toFinFlag = ⟨r, Flag_r_0_0_idx⟩}. The only datum
left is \btt{idx}, the canonical index of $K_r$ among the empty-typed
$r$-vertex flags ($K_3\!\to\!3$, $K_4\!\to\!10$, $K_5\!\to\!33$).

\paragraph{The two forbid loaders --- was the easy-to-miss Step 3, now nothing.}
This is the step Section~\ref{sec:newforbid} singled out as manual and easy to
forget. Both \texttt{load\_forbid\_density\_theorems} and
\texttt{load\_forbid\_mul\_theorems} used a \texttt{match} on the tag with one
copy-pasted \btt{| "K3" =>} / \btt{"K4"} / \btt{"K5"} branch apiece (each baking
in a lemma name and a literal \btt{simp [Flag_r_0_0_idx]}). Those branches are
gone. A helper \btt{parseCompleteGraphTag} recovers $r$ from a \btt{"K<r>"} tag,
and --- the piece that removes the last hard-coded index ---
\texttt{load\_forbid\_density\_theorems} now reads the forbidden flag's
representative \emph{off} the \btt{K{r}_toFinFlag_eq} lemma (via
\btt{forbidFlagIdentOfToFinFlagEq}) instead of from a literal index. The
multiplication loader never needed the index, so it just refers to \btt{K{r}}
by name.

\paragraph{Net effect.} Adding a new clique forbid $K_r$ no longer touches five
files: run the Python step with \btt{--forbid-Kn r}, and add one line
\btt{generate_complete_graph r <idx>} to \texttt{CommonGraphs.lean}. The loaders
need no edit --- Step~3 of Section~\ref{sec:newforbid} disappears for complete
graphs, leaving only the one genuine datum (the index $i_{K_r}$). The refactor
was checked by rebuilding the $K_3$/$K_4$/$K_5$ certificates (Mantel, K4turan,
K5turan) green.

\paragraph{Still bounded.} The empty-typed flags exist only up to $n=5$
(\btt{generate_empty_typed_flags 6} is commented out in
\texttt{FlagDef.lean}), so $K_n$ forbids are available through $n=5$ today; $K_6$
and up wait on enabling those flags.

\section{What is available now}

\paragraph{Sizes.} Empty-typed flags up to $n = 6$; typed flags up to
$n = 5$ (the type/size pairs listed in \texttt{FlagDef.lean}).

\paragraph{Forbidden graphs.} $K_3$, $K_4$, $K_5$ are wired end to end.
Adding another one is the procedure of Section~\ref{sec:newforbid} (five
files, three steps); the certificate translator needs no change. Non-$K_n$
forbids ($C_4$, $P_4$, \ldots) have not been run end to end yet.

\paragraph{Ported certificates.} Six build with no \texttt{sorry}:

\begin{center}
\begin{tabular}{lcccc}
\toprule
Certificate & Forbid & $N$ & blocks & bound \\
\midrule
Mantel        & $K_3$ & 3 & 1 & $1/2$    \\
K3forbidP3    & $K_3$ & 3 & 1 & $3/4$    \\
K3forbidC4    & $K_3$ & 4 & 2 & $3/8$    \\
K4turan       & $K_4$ & 4 & 2 & $2/3$    \\
ErdosPentagon & $K_3$ & 5 & 3 & $24/625$ \\
K5turan       & $K_5$ & 5 & 4 & $3/4$    \\
\bottomrule
\end{tabular}
\end{center}

\paragraph{Manual steps that remain.} Producing the density JSON is still
a Python step (Group~1 flag definitions are not). A new forbidden graph
needs the procedure of Section~\ref{sec:newforbid}. The branch
that lifts the objective to host size ends in a \texttt{ring\_nf} whose
success has not been characterized in general, though it has closed every
case so far.

\end{document}
